% ---------------------------------------------------------------------------
%  T1 — Actividad 1 (Ética, parte 1)
%  Presentación MINIMALISTA: "Ética en la Publicación Científica"
%  Grupo 1 de Cátedra — Metodologías de Investigación Aplicada (MIA)
%  Tema: usach (normas gráficas MNG 2024-2S) + iconos FontAwesome v4
%  Compilar: latexmk -pdf -outdir=build main-minimalista.tex
% ---------------------------------------------------------------------------
\documentclass[11pt,aspectratio=169,xcolor=table]{beamer}
\pdfmapfile{+fontawesome.map}

\IfFileExists{spanish.ldf}{%
  \usepackage[spanish,es-noquoting,es-nodecimaldot]{babel}%
}{%
  \usepackage[T1]{fontenc}
  \renewcommand{\figurename}{Figura}
  \renewcommand{\tablename}{Tabla}
}
\usepackage[utf8]{inputenc}
\usepackage{amsmath,amssymb}
\usepackage{booktabs,tabularx}
\usepackage{microtype}
\usepackage{fontawesome}
\usepackage{tikz}
\graphicspath{{images/}}

\usetheme{usach}

% --- Agenda con icono por seccion (mismos colores que las portadillas) -------
\setbeamertemplate{section in toc}{%
  \leavevmode
  \hbox to 0.62cm{{\color{usachrojo}\usachtitular\inserttocsectionnumber}\hss}%
  {\ifcase\inserttocsectionnumber\relax
     \color{usachverde}\faBook
   \or\color{usachverde}\faBook
   \or\color{usachnaranjo}\faFlask
   \or\color{usachazul}\faIndustry
   \or\color{usachpurpura}\faBalanceScale
   \fi}%
  \hskip0.45em
  \inserttocsection\par}

% --- Helpers minimalistas ----------------------------------------------------
% Viñeta con icono: \iconitem[color]{icono}{Titulo}{detalle}
\newcommand{\iconitem}[4][usachrojo]{%
  \item[\textcolor{#1}{#2}] \textbf{#3} — #4}

% Dato destacado con icono: \usachdatoi[color]{icono}{cifra}{detalle}
\newcommand{\usachdatoi}[4][usachrojo]{%
  \begin{minipage}[t]{\linewidth}
    {\color{#1}\fontsize{15}{17}\selectfont #2\hspace{0.45em}%
     \raisebox{0.45em}{\color{#1}\rule{0.85cm}{1.6pt}}\par}
    \vskip0.35em
    {\color{usachtinta}\usachcifra #3\par}
    \vskip0.10em
    {\color{usachgris70}\usachcuerpo\fontsize{8.5}{10.5}\selectfont #4\par}
  \end{minipage}}

% Retrato circular con nombre y rol: \usachfoto[color]{archivo}{Nombre}{Rol/afiliacion}
% Si images/<archivo> no existe todavia, dibuja un avatar generico (no rompe la compilacion).
\newcommand{\usachfoto}[4][usachrojo]{%
  \begin{minipage}[t]{2.05cm}\centering
    \IfFileExists{images/#2}{%
      \begin{tikzpicture}
        \begin{scope}
          \clip (0,0) circle (0.8cm);
          \node[inner sep=0pt] at (0,0)
            {\includegraphics[width=1.8cm,height=1.8cm,keepaspectratio=false]{images/#2}};
        \end{scope}
        \draw[#1,line width=1.1pt] (0,0) circle (0.8cm);
      \end{tikzpicture}%
    }{%
      \begin{tikzpicture}
        \draw[fill=#1!12,draw=#1,line width=1.1pt] (0,0) circle (0.8cm);
        \node[inner sep=0pt] at (0,0) {\textcolor{#1}{\fontsize{20}{20}\selectfont\faUser}};
      \end{tikzpicture}%
    }
    \vskip0.22em
    {\usachcuerpo\fontsize{7.5}{9}\selectfont\color{usachtinta}\textbf{#3}\par}
    {\usachcuerpo\fontsize{6.5}{8}\selectfont\color{usachgris70}#4\par}
  \end{minipage}}

% Logo institucional con alto maximo: \usachlogo[alto]{archivo}{Texto alternativo}
% Si images/<archivo> no existe todavia, dibuja una caja con el nombre (no rompe la compilacion).
\newcommand{\usachlogo}[3][1.1cm]{%
  \IfFileExists{images/#2}{%
    \includegraphics[height=#1]{images/#2}%
  }{%
    \tikz\node[draw=usachgris20,fill=usachgris10,rounded corners=2pt,
      minimum width=2.4cm,minimum height=#1,align=center,inner sep=3pt]
      {\usachcuerpo\fontsize{7}{8}\selectfont\color{usachgris70}#3};
  }}

% Portadilla de seccion minimalista: el tema inyecta \usachpaginaseccion
% automaticamente en cada \section (AtBeginSection); la redefinimos sin el
% TOC lateral para mantener la estetica minimalista.
\renewcommand{\usachpaginaseccion}{%
  \begin{usachlaminaoscura}[c]
    \begin{columns}[T]
      \begin{column}{0.14\textwidth}
        {\color{usachverde}\fontsize{30}{34}\selectfont
          \ifcase\insertsectionnumber\relax\faBook
            \or\faBook\or\faFlask\or\faIndustry\or\faBalanceScale\fi}
      \end{column}
      \begin{column}{0.82\textwidth}
        \usachpleca[white]{%
          {\usachcuerpo\fontsize{9}{11}\selectfont\color{usachgris40}
            Secci\'on \insertsectionnumber\par}
          \vskip0.5em
          {\usachtitular\fontsize{24}{27}\selectfont\color{white}\insertsectionhead\par}}
      \end{column}
    \end{columns}
  \end{usachlaminaoscura}}

% --- Metadatos --------------------------------------------------------------
\title[Ética en la publicación científica]{Ética en la Publicación Científica}
\subtitle{De la manipulación individual de datos a la industrialización del fraude}
\author[G.~Ahumada · D.~Fernández · I.~Fernández · P.~Figueroa]{%
  Gonzalo Ahumada · Diego Fernández · Iván Fernández · Pablo Figueroa}
\date{Septiembre de 2026}

\usachfacultad{Facultad de Ingeniería}
\usachdepartamento{Departamento de Ingeniería Informática}
\usachprograma{Magíster en Ingeniería Informática}
\usachasignatura{Metodologías de Investigación Aplicada}
\usachprofesor{Prof. Roberto Ibáñez González}

\institute[USACH]{Universidad de Santiago de Chile}

\begin{document}

\usachportada
\note{Somos el Grupo 1 de cátedra. Presentamos dos casos reales de ética en
la publicación: Jan Hendrik Schön y las retracciones masivas de la ACM.}

% --- Agenda ------------------------------------------------------------------
\begin{frame}{Contenido}
  \tableofcontents
\end{frame}

% =============================================================================
\section{Marco normativo}
% =============================================================================

\begin{frame}{La publicación como producto}{Tríada FFP}
  La publicación es el producto formal y acreditable de la investigación.
  \begin{itemize}\usachaire
    \iconitem[usachrojo]{\faFlask}{Fabricación}{inventar datos o resultados}
    \iconitem[usachnaranjo]{\faPencilSquareO}{Falsificación}{manipular datos reales}
    \iconitem[usachpurpura]{\faCopy}{Plagio}{apropiarse sin atribución}
  \end{itemize}
  \usachpleca[usachgris20]{Foco en la publicación como producto; se excluyen
  las faltas contra personas.}
  \note{Tríada FFP según \emph{On Being a Scientist} y la Declaración de
  Singapur.}
\end{frame}

\begin{frame}{Singapur y presiones}{Deberes y riesgos}
  \begin{columns}[T,onlytextwidth]
    \begin{column}{0.48\textwidth}
      \textbf{Deberes de Singapur}
      \begin{itemize}\usachaire
        \iconitem[usachverde]{\faDatabase}{Resp.~4}{custodia de datos}
        \iconitem[usachverde]{\faUsers}{Resp.~6}{autoría responsable}
        \iconitem[usachverde]{\faEye}{Resp.~8}{arbitraje riguroso}
      \end{itemize}
    \end{column}
    \begin{column}{0.48\textwidth}
      \textbf{Presiones y modalidades}
      \begin{itemize}\usachaire
        \iconitem[usachrojo]{\faLineChart}{\emph{Publish or perish}}{volumen como incentivo}
        \iconitem[usachrojo]{\faUserSecret}{Artesanal}{falsificador solitario}
        \iconitem[usachrojo]{\faIndustry}{Industrializado}{\emph{paper mills}}
        \iconitem[usachrojo]{\faGlobe}{Vigilancia}{Retraction Watch y PubPeer}
      \end{itemize}
    \end{column}
  \end{columns}
  \note{Singapur: custodia de datos, autoría solidaria y revisión por pares
  rigurosa. El fraude evolucionó del individuo a la industria.}
\end{frame}

% =============================================================================
\section{Caso 1: Jan Hendrik Schön}
% =============================================================================

\begin{frame}{El «prodigio» de Bell Labs}
  {\def\usachcifra{\usachtitular\fontsize{22}{24}\selectfont}
  \usachdatos{%
    \begin{minipage}[t]{0.23\linewidth}\raggedright\hyphenpenalty=10000\exhyphenpenalty=10000 \usachdatoi[usachverde]{\faCalendarO}{8 días}{entre cada paper}\end{minipage}\hfill
    \begin{minipage}[t]{0.23\linewidth}\raggedright\hyphenpenalty=10000\exhyphenpenalty=10000 \usachdatoi[usachazul]{\faFileTextO}{4 años}{publicando (1998–2001)}\end{minipage}\hfill
    \begin{minipage}[t]{0.23\linewidth}\raggedright\hyphenpenalty=10000\exhyphenpenalty=10000 \usachdatoi[usachnaranjo]{\faTrophy}{<35 años}{candidato al Nobel}\end{minipage}\hfill
    \begin{minipage}[t]{0.23\linewidth}\raggedright\hyphenpenalty=10000\exhyphenpenalty=10000 \usachdatoi[usachpurpura]{\faUsers}{20}{coautores secundarios}\end{minipage}}}
  \begin{itemize}\usachaire
    \iconitem[usachrojo]{\faFlask}{Bell Labs + Konstanz}{física de la materia condensada}
    \iconitem[usachrojo]{\faBook}{\emph{Science}, \emph{Nature} y APL}{transistores monomoleculares}
    \iconitem[usachrojo]{\faUserSecret}{Efecto halo}{el prestigio silenció el escepticismo}
  \end{itemize}
  \vskip0.35em
  \begin{center}
    \usachfoto[usachrojo]{schon.jpg}{Jan H. Schön}{Investigador principal}%
    \hspace{1.6em}%
    \usachfoto[usachrojo]{batlogg.jpg}{Bertram Batlogg}{Coautor sénior}%
    \hspace{2.4em}%
    \raisebox{0.5cm}{\usachlogo[0.8cm]{logo-bell-labs.png}{Bell Labs}}%
    \hspace{1.2em}%
    \raisebox{0.5cm}{\usachlogo[0.8cm]{logo-konstanz.png}{Universität Konstanz}}
  \end{center}
  \note{Un paper cada ocho días y candidato al Nobel antes de los 35. El
  prestigio de Bell Labs creó una burbuja de confianza ciega.}
\end{frame}

\begin{frame}{Fabricar, falsificar y borrar}{Transgresiones}
  \begin{columns}[T,onlytextwidth]
    \begin{column}{0.48\textwidth}
      \textbf{Según «On Being a Scientist»}
      \begin{itemize}\usachaire
        \iconitem[usachrojo]{\faLineChart}{Curvas ideales}{fórmulas + ruido artificial}
        \iconitem[usachrojo]{\faTrashO}{Datos «borrados»}{sin bitácoras; chips a la basura}
        \iconitem[usachgris70]{\faUsers}{Coautores sénior}{reaccionaron bien; ¿debieron ser proactivos?}
      \end{itemize}
    \end{column}
    \begin{column}{0.48\textwidth}
      \textbf{Contra Singapur}
      \begin{itemize}\usachaire
        \iconitem[usachnaranjo]{\faBan}{Honestidad}{adulteración por ventaja}
        \iconitem[usachnaranjo]{\faDatabase}{Resp.~4}{sin registros reproducibles}
        \iconitem[usachgris70]{\faUsers}{Resp.~6}{responsabilidad de coautoría en debate}
      \end{itemize}
    \end{column}
  \end{columns}
  \note{No hacía ciencia experimental: graficaba funciones ideales con ruido
  simulado y borró toda evidencia.}
\end{frame}

\begin{frame}{El ruido gemelo}{Revelación y sanciones}
  \begin{columns}[T,onlytextwidth]
    \begin{column}{0.48\textwidth}
      \textbf{La revelación (2002)}
      \begin{itemize}\usachaire
        \iconitem[usachazul]{\faFlask}{Réplicas fallidas}{en múltiples laboratorios}
        \iconitem[usachazul]{\faClone}{Ruido idéntico}{píxel por píxel: imposible}
        \iconitem[usachazul]{\faSearch}{Sohn y McEuen}{cotejaron dos papers}
      \end{itemize}
    \end{column}
    \begin{column}{0.48\textwidth}
      \begin{columns}[c,onlytextwidth]
        \begin{column}{0.20\textwidth}
          \usachfoto[usachrojo]{beasley.jpg}{M. Beasley}{Stanford}
        \end{column}
        \begin{column}{0.78\textwidth}
          \hspace{0.4cm}\textbf{Comité Beasley}
        \end{column}
      \end{columns}
      \begin{itemize}\usachaire
        \iconitem[usachrojo]{\faExclamationTriangle}{Fraude confirmado}{en al menos 16 papers}
        \iconitem[usachrojo]{\faBan}{Retractaciones}{10 en \emph{Science}, 7 en \emph{Nature}}
        \iconitem[usachrojo]{\faGraduationCap}{Doctorado revocado}{y despido de Bell Labs}
      \end{itemize}
    \end{column}
  \end{columns}
  \note{Lo detectó la comunidad al no poder replicar: dos experimentos
  distintos compartían el mismo ruido, copiado píxel a píxel.}
\end{frame}

% =============================================================================
\section{Caso 2: escándalo ACM}
% =============================================================================

\begin{frame}{ACM: \emph{paper mills}}{Mercantilización de la autoría}
  {\def\usachcifra{\usachtitular\fontsize{19}{22}\selectfont}
  \usachdatos{%
    \begin{minipage}[t]{0.40\linewidth}\raggedright\hyphenpenalty=10000\exhyphenpenalty=10000 \usachdatoi[usachrojo]{\faBan}{>320}{artículos retractados}\end{minipage}\hfill
    \begin{minipage}[t]{0.30\linewidth}\centering\vspace{0.3em}\usachlogo[1.4cm]{logo-acm.jpg}{ACM}\end{minipage}}}
  \begin{itemize}\usachaire
    \iconitem[usachrojo]{\faIndustry}{ACM e ICPS}{actas de ICIMTech afectadas}
    \iconitem[usachrojo]{\faMoney}{Autorías subastadas}{incentivos por publicar en Scopus}
  \end{itemize}
  \note{La ACM, la sociedad de computación más prestigiosa, sufrió fábricas
  de artículos que vendían autorías.}
\end{frame}

\begin{frame}{«Tortured phrases»}{Cómo burlaron los filtros antiplagio}
  \begin{table}
    \centering\footnotesize
    \begin{tabularx}{\textwidth}{@{}X X X@{}}
      \toprule
      \rowcolor{usachgris}
      \usachcab{Término real} & \usachcab{Tortured phrase} & \usachcab{Contexto} \\
      \midrule
      Artificial Intelligence & Counterfeit Consciousness & Modelos de ML / IA \\
      Deep Learning & Profound Learning & Redes neuronales \\
      Cloud Computing & Haze Computing & Arquitecturas en la nube \\
      Random Forest & Arbitrary Timberland & Clasificación supervisada \\
      \bottomrule
    \end{tabularx}
    \caption{Ejemplos de parafraseo inverso algorítmico.}
  \end{table}
  \usachfuente{Cabanac, Labbé y Magazinov (2021).}
  \begin{itemize}\usachaire
    \iconitem[usachrojo]{\faRandom}{\emph{Text spinners}}{burlan Turnitin/iThenticate}
    \iconitem[usachrojo]{\faUsers}{Revisores ficticios}{aprobación exprés, sin revisión real}
  \end{itemize}
  \note{Los parafraseadores no entienden la jerga: «Artificial Intelligence»
  se volvió «Counterfeit Consciousness» y aun así se aprobaron.}
\end{frame}

\begin{frame}{Minería de texto y sanciones}{La computación desenmascara el fraude}
  \begin{columns}[T,onlytextwidth]
    \begin{column}{0.48\textwidth}
      \textbf{Detectives digitales}
      \vskip0.3em
      \begin{columns}[T,onlytextwidth]
        \begin{column}{0.31\textwidth}
          \usachfoto[usachazul]{cabanac.jpg}{Cabanac}{U. Toulouse III}
        \end{column}
        \begin{column}{0.31\textwidth}
          \usachfoto[usachazul]{labbe.jpg}{Labbé}{U. Grenoble Alpes}
        \end{column}
        \begin{column}{0.31\textwidth}
          \usachfoto[usachazul]{magazinov.jpg}{Magazinov}{Yandex / Skoltech}
        \end{column}
      \end{columns}
      \vskip0.4em
      \begin{itemize}\usachaire
        \iconitem[usachazul]{\faCode}{\emph{Problematic Paper Screener}}{herramienta de minería de texto}
        \iconitem[usachazul]{\faNewspaperO}{PubPeer y \emph{Nature}}{presión pública}
      \end{itemize}
    \end{column}
    \begin{column}{0.48\textwidth}
      \textbf{Consecuencias}
      \begin{itemize}\usachaire
        \iconitem[usachrojo]{\faBan}{Abril 2022}{retractación masiva (>320)}
        \iconitem[usachrojo]{\faTimesCircleO}{Veto}{duración variable, sin publicar}
        \iconitem[usachrojo]{\faDatabase}{\emph{Data poisoning}}{riesgo teórico para IA futura}
      \end{itemize}
    \end{column}
  \end{columns}
  \note{Cabanac, Labbé y Magazinov programaron un detector de frases
  torturadas. La ACM retractó más de 320 artículos.}
\end{frame}

% =============================================================================
\section{Comparación y lecciones}
% =============================================================================

\begin{frame}{Schön vs. ACM}{Dos décadas de fraude}
  \begin{table}
    \centering\footnotesize
    \begin{tabularx}{\textwidth}{@{}X X X@{}}
      \toprule
      \rowcolor{usachgris}
      \usachcab{Dimensión} & \usachcab{Schön (física)} & \usachcab{ACM (informática)} \\
      \midrule
      Naturaleza & Artesanal e individual & Industrializada y sistémica \\
      Mecanismo & Curvas ideales + ruido simulado & \emph{Text spinners} antiplagio \\
      Falla & Confianza ciega en el prestigio & Control editorial tercerizado \\
      Detección & Réplica fallida + ruido idéntico & Minería de texto y PubPeer \\
      Norma rota & FFP + Resp.~1, 3, 4; Resp.~6 en debate & Resp.~1, 6 (autoría) y Resp.~8 (arbitraje) \\
      \bottomrule
    \end{tabularx}
    \caption{Del individuo a la red transnacional.}
  \end{table}
  \usachfuente{Beasley et al. (2002), Cabanac et al. (2021) y Retraction Watch (2022).}
  \note{A Schön lo delató la imposibilidad de replicar; a la ACM, la propia
  computación mediante minería de texto.}
\end{frame}

\usachdestacado{Del fraude artesanal al industrializado, la integridad
científica es el único activo que valida nuestra labor.}

\begin{frame}{Lecciones para el posgrado}{Informática aplicada}
  \begin{columns}[T,onlytextwidth]
    \begin{column}{0.48\textwidth}
      \textbf{Pilares innegociables}
      \begin{itemize}\usachaire
        \iconitem[usachverde]{\faDatabase}{Datos y código abiertos}{«perdí los datos» confiesa mala praxis}
        \iconitem[usachverde]{\faUsers}{Autoría solidaria}{no existe autoría honoraria inocua}
        \iconitem[usachverde]{\faEye}{Arbitraje moral}{sin bots ni acuerdos recíprocos}
      \end{itemize}
    \end{column}
    \begin{column}{0.48\textwidth}
      \textbf{Desafíos de la era IA}
      \begin{itemize}\usachaire
        \iconitem[usachrojo]{\faCogs}{LLM como asistencia}{nunca pseudo-conocimiento}
        \iconitem[usachrojo]{\faGlobe}{Vigilancia post-publicación}{PubPeer y Retraction Watch}
        \iconitem[usachrojo]{\faBalanceScale}{Rigor antes que volumen}{superar el \emph{publish or perish}}
      \end{itemize}
    \end{column}
  \end{columns}
  \note{Reproducibilidad, autoría solidaria y guardianes de la calidad
  epistemológica de los datos.}
\end{frame}

% --- Fuentes -----------------------------------------------------------------
\begin{frame}{Fuentes}
  \begin{itemize}\usachaire
    \iconitem[usachverde]{\faBook}{National Academy of Sciences}{\emph{On Being a Scientist} (3.ª ed.), p.~35}
    \iconitem[usachverde]{\faBook}{WCRI}{\emph{Declaración de Singapur} (2010)}
    \iconitem[usachverde]{\faBook}{Retraction Watch}{base de datos de retracciones}
    \iconitem[usachverde]{\faBook}{Cabanac, Labbé y Magazinov (2021)}{\emph{Tortured phrases}}
    \iconitem[usachverde]{\faBook}{Beall's List}{revistas predatorias}
  \end{itemize}
\end{frame}

\usachcierre[¿Preguntas y comentarios?]{Gracias}

\end{document}